\documentclass[12pt]{article}

% Preamble and title page follow the TU-Sofia thesis template (Bulgarian layout,
% same faculty heading and typography).

\usepackage{cmap}
\usepackage[T2A]{fontenc}
\usepackage[utf8]{inputenc}
\usepackage[bulgarian]{babel}

% --- Typography ---
\usepackage{tempora}
\usepackage[top=2.5cm, bottom=2.5cm, left=3.0cm, right=2.0cm, a4paper]{geometry}
\usepackage{setspace}
\setstretch{1.5}
\usepackage[none]{hyphenat}
\usepackage[expansion=false]{microtype}
\emergencystretch=1.5em
\sloppy
\usepackage{indentfirst}
\setlength{\parindent}{1.25cm}
\setlength{\parskip}{0pt}

% --- Page numbering (bottom right) ---
\usepackage{fancyhdr}
\pagestyle{fancy}
\fancyhf{}
\renewcommand{\headrulewidth}{0pt}
\fancyfoot[R]{\thepage}

% --- Footnotes ---
\usepackage[hang,flushmargin]{footmisc}
\usepackage{ragged2e}
\renewcommand{\footnotelayout}{\fontsize{10}{12}\selectfont\setstretch{1.0}\justifying}

% --- Sections ---
\usepackage{titlesec}
\renewcommand{\thesection}{\Roman{section}}
\renewcommand{\thesubsection}{\arabic{section}.\arabic{subsection}}

\titleformat{\section}[block]
  {\fontsize{14}{16.8}\bfseries\centering}
  {\thesection. }
  {0pt}
  {\MakeUppercase}

\titleformat{\subsection}[block]
  {\fontsize{14}{16.8}\bfseries}
  {\thesubsection. }
  {0pt}
  {}
\titlespacing*{\subsection}{1.25cm}{*3}{*2}

% --- Table of contents ---
\usepackage{etoolbox}
\makeatletter
\patchcmd{\l@section}{#1}{\MakeUppercase{#1}}{}{}
\makeatother
\usepackage[titles]{tocloft}
\renewcommand{\cftsecleader}{\cftdotfill{\cftdotsep}}
\renewcommand{\cftsubsecleader}{\cftdotfill{\cftdotsep}}
\renewcommand{\cftsecfont}{\fontsize{12}{14}\bfseries}
\renewcommand{\cftsecpagefont}{\fontsize{12}{14}\bfseries}
\renewcommand{\cftsecaftersnum}{.}
\renewcommand{\cftsubsecaftersnum}{.}

% --- Captions ---
\usepackage{caption}
\DeclareCaptionLabelFormat{tablelabel}{Таблица #2}
\captionsetup[table]{
    labelformat=tablelabel,
    labelsep=period,
    font={bf,small},
    justification=centering,
    singlelinecheck=true
}

% --- Other ---
\usepackage{amsmath}
\usepackage{graphicx}
\usepackage{longtable}
\usepackage{needspace}
\usepackage{listings}
\lstset{basicstyle=\ttfamily\footnotesize,breaklines=true,columns=fullflexible,keepspaces=true,frame=single}
\usepackage{float}
\usepackage{booktabs}
\usepackage{array}
\usepackage[hidelinks]{hyperref}

\AtBeginEnvironment{lstlisting}{\setstretch{1}}
\begin{document}
\hypersetup{pageanchor=false}
\begin{titlepage}
\begin{center}
{\fontsize{18}{21.6}\selectfont\bfseries ТЕХНИЧЕСКИ УНИВЕРСИТЕТ -- СОФИЯ\par}
\vspace{0.5cm}
{\fontsize{16}{19.2}\selectfont\bfseries ФАКУЛТЕТ „КОМПЮТЪРНИ СИСТЕМИ И ТЕХНОЛОГИИ“\par}
\vspace{0.6cm}
\includegraphics[width=0.28\linewidth]{images/tu_logo.png}
\vspace{0.8cm}

{\fontsize{20}{24}\selectfont\bfseries ДИАГНОСТИКА НА PPP С CHAP УДОСТОВЕРЯВАНЕ\par}
\vspace{0.8cm}
{\fontsize{14}{16.8}\selectfont\bfseries ЛАБОРАТОРЕН ПРОТОКОЛ\par}
\vspace{0.3cm}
{\fontsize{13}{15.6}\selectfont Упражнение 2\par Cisco Packet Tracer\par}
\vfill
\begin{flushleft}
{\fontsize{14}{16.8}\selectfont\bfseries Изготвил:\par}
Алекс Цветанов\par
Факултетен № 371326005\par
\vspace{0.5cm}
{\fontsize{14}{16.8}\selectfont\bfseries Преподавател:\par}
д-р Йонка Ганчева\par
\end{flushleft}
\vfill
София, 26 октомври 2026 г.
\end{center}
\end{titlepage}
\hypersetup{pageanchor=true}
\pagenumbering{roman}
\tableofcontents
\clearpage
\pagenumbering{arabic}

\section{Цел и постановка}
Задачата „Packet Tracer -- Troubleshooting PPP with Authentication“ изисква възстановяване на свързаността между три маршрутизатора, две локални мрежи и външен уеб сървър. Диагностиката обхваща физическото, каналното и мрежовото ниво. Всички серийни връзки трябва да използват PPP и CHAP с парола \texttt{cisco}.

Упражнението е изпълнено в Cisco Packet Tracer 9.0.1. Началната оценка е 0/100, а след корекциите автоматичната проверка отчита \textbf{100/100}. Допълнително са изпълнени проверки с \texttt{ping} и \texttt{tracert} от крайните станции.

\section{Топология и адресиране}
Маршрутизаторите R1, R2 и R3 са свързани в триъгълник. PC1 достига R1 през S1, а PC3 достига R3 през S3. R2 е свързан с ISP, зад който се намира Web Server.

\begin{center}\small
\begin{longtable}{llll}
\toprule Устройство & Интерфейс & IPv4 адрес & Маска \\\midrule
R1 & G0/1 & 10.0.0.1 & /25 \\
R1 & S0/0/0 & 172.16.0.1 & /30 \\
R1 & S0/0/1 & 172.16.0.9 & /30 \\
R2 & G0/1 & 209.165.200.161 & /27 \\
R2 & S0/0/0 & 172.16.0.2 & /30 \\
R2 & S0/0/1 & 172.16.0.5 & /30 \\
R3 & G0/1 & 10.0.0.129 & /25 \\
R3 & S0/0/0 & 172.16.0.10 & /30 \\
R3 & S0/0/1 & 172.16.0.6 & /30 \\
ISP & G0/1 & 209.165.200.162 & /27 \\
PC1 & NIC & 10.0.0.10 & /25 \\
PC3 & NIC & 10.0.0.139 & /25 \\
Web Server & NIC & 209.165.200.2 & /30 \\\bottomrule
\end{longtable}\end{center}
Шлюзовете на PC1, PC3 и Web Server са съответно 10.0.0.1, 10.0.0.129 и 209.165.200.1. Маските /25, /27 и /30 съответстват на 255.255.255.128, 255.255.255.224 и 255.255.255.252.

\section{Диагностика и установени грешки}
\subsection{Физическо ниво}
Проверени са портовете на кабелите и административното състояние на интерфейсите. Кабелните връзки съответстват на топологията и не се налага преместване. LAN кабелът на R1 е включен в G0/1, въпреки че първоначално адресът е зададен на изключения G0/0.

DCE краищата са R1 S0/0/0, R2 S0/0/1 и R3 S0/0/0. Наличните тактови честоти са съответно 64000, 2000000 и 64000 bit/s. Те са запазени. Неактивните необходими интерфейси са включени с \texttt{no shutdown}.

\subsection{Канално ниво и CHAP}
В началната конфигурация част от серийните портове използват HDLC, а други PPP с PAP. На всички шест серийни интерфейса е осигурено \texttt{encapsulation ppp} и \texttt{ppp authentication chap}. Съществуващите правилни настройки са запазени.

На R1 липсва потребител R3; на R3 липсва R2. На R2 е въведено грешно име R11 вместо R1, а паролата за R3 е \texttt{class}. Записите са поправени така, че всеки маршрутизатор разполага с имената на двата си непосредствени съседи и общата парола \texttt{cisco}.

\subsection{Мрежово ниво}
На R1 адресът 10.0.0.1/25 е преместен от G0/0 на G0/1. Маската на S0/0/0 е коригирана от /29 на /30. На R2 е добавен липсващият адрес 209.165.200.161/27 на G0/1. На R3 маската на G0/1 е променена от /24 на /25, а на S0/0/1 е зададен липсващият адрес 172.16.0.6/30.

След възстановяването на интерфейсите OSPF съседствата достигат състояние FULL. Първите проверки към сървъра обаче връщат „Destination host unreachable“ от LAN шлюзовете. В R1 и R3 липсва маршрут по подразбиране, въпреки наличните статичен маршрут и \texttt{default-information originate} в R2. В OSPF базата е наблюдавана външна обява за 0.0.0.0 с Advertising Router 0.0.0.0.

Рестартирането на OSPF процеса в R2 възстановява съседствата, но само по себе си не решава проблема. Премахването и повторното задаване на \texttt{default-information originate} обновява обявата. След тази операция крайните станции достигат сървъра. Това е допълнителна проверка на работата на мрежата извън автоматичната оценка.

\section{Приложени команди}
Следващите списъци съдържат приложените корекции, а не пълни конфигурации. Конфигурациите са записани с \texttt{write memory}.
\subsection{R1}
\begin{lstlisting}
configure terminal
username R3 password cisco
interface g0/0
 no ip address
interface g0/1
 ip address 10.0.0.1 255.255.255.128
 no shutdown
interface s0/0/0
 ip address 172.16.0.1 255.255.255.252
 encapsulation ppp
 ppp authentication chap
 no shutdown
interface s0/0/1
 ppp authentication chap
 no shutdown
end
write memory
\end{lstlisting}
\subsection{R2}
\begin{lstlisting}
configure terminal
no username R11
username R1 password cisco
username R3 password cisco
interface g0/1
 ip address 209.165.200.161 255.255.255.224
interface s0/0/0
 no shutdown
interface s0/0/1
 encapsulation ppp
 ppp authentication chap
 no shutdown
end
write memory
clear ip ospf process
yes
configure terminal
router ospf 1
 no default-information originate
 default-information originate
end
write memory
\end{lstlisting}
\subsection{R3}
\begin{lstlisting}
configure terminal
username R2 password cisco
interface g0/1
 ip address 10.0.0.129 255.255.255.128
 no shutdown
interface s0/0/0
 encapsulation ppp
 ppp authentication chap
 no shutdown
interface s0/0/1
 ip address 172.16.0.6 255.255.255.252
 ppp authentication chap
 no shutdown
end
write memory
\end{lstlisting}

\section{Проверки и резултати}
Автоматичната проверка отчита \textbf{Completion: 100/100} и съобщение за успешно завършване на упражнението. Серийните връзки са активни. На R3 са наблюдавани G0/1, S0/0/0 и S0/0/1 в състояние \texttt{up/up} и две OSPF съседства в състояние FULL.

Първото изпълнение на ping след обновяването на OSPF обявата дава по три отговора от четири и на двата компютъра. Повторните проверки приключват без загуби:
\begin{center}
\begin{tabular}{lrrrr}
\toprule Източник & Получени & Загуби & min/avg/max, ms \\\midrule
PC1 & 4/4 & 0\% & 2/32/64 \\
PC3 & 4/4 & 0\% & 40/57/70 \\\bottomrule
\end{tabular}
\end{center}
И при двата теста целта е 209.165.200.2, полезният товар е 32 bytes, а полученият TTL е 125. Стойностите са измервания в симулатора.

\subsection{Проследяване от PC1}
\begin{lstlisting}
C:\>tracert 209.165.200.2
1   0 ms    0 ms    0 ms    10.0.0.1
2  28 ms   22 ms    0 ms    172.16.0.2
3  10 ms   30 ms    9 ms    209.165.200.162
4  11 ms    5 ms    0 ms    209.165.200.2
Trace complete.
\end{lstlisting}
Пътят преминава през R1, R2 и ISP и достига сървъра за четири стъпки.

\subsection{Проследяване от PC3}
\begin{lstlisting}
C:\>tracert 209.165.200.2
1   0 ms    0 ms    0 ms    10.0.0.129
2  10 ms    5 ms    8 ms    172.16.0.5
3   5 ms    0 ms    9 ms    209.165.200.162
4   0 ms    0 ms    9 ms    209.165.200.2
Trace complete.
\end{lstlisting}
Пътят преминава през R3, R2 и ISP. Двата успешни теста потвърждават достижимостта на външния сървър от двете локални мрежи.

\Needspace{12\baselineskip}
\section{Заключение}
Отстранени са административно изключени интерфейси, несъответстващи PPP настройки, грешки в CHAP записите и IPv4 адресирането. След обновяване на OSPF обявата за маршрута по подразбиране е възстановена и свързаността към външния сървър. Крайният резултат е 100/100, а повторните ping проверки са с 0\% загуби.

Към протокола са приложени началният и завършеният PKA файл, командите по устройства и текстов препис на наблюдаваните резултати. Командните файлове съдържат корекциите, а не пълните running-config конфигурации.
\end{document}
